% \documentstyle[twocolumn, 10pt]{article}
\documentclass[10pt]{article}
\usepackage{latexsym}

        \textwidth 6.5in
        \textheight 9in
\evensidemargin .05in
        \oddsidemargin .05in
        \topmargin 0in  
        \headsep 0in    
        \headheight 0in
        \footskip .5in
    \title{Formula Sheet For Math 152}
    \author{}
    \date{}
    \begin{document}
\twocolumn
%\maketitle


\noindent {\large \bf Numerarical Information}


$$\begin{array}{cccccc}
\theta&0&\frac{\pi}{6}&\frac{\pi}{4}&\frac{\pi}{3}&\frac{\pi}{2}\\
\sin \theta&0&\frac{1}{2}&\frac{\sqrt{2}}{2}&\frac{\sqrt{3}}{2}&1\\
\cos \theta&1&\frac{\sqrt{3}}{2}&\frac{\sqrt{2}}{2}&\frac{1}{2}&0
\end{array}$$

\noindent {\large \bf Differentiation Formulas}

$$\frac{d}{dx}(\sin^{-1} x)=\frac{1}{\sqrt{1-x^{2}}}$$
$$\frac{d}{dx}(\tan^{-1} x)=\frac{1}{1+x^{2}}$$

\noindent {\large \bf Integration Formulas}

\begin{eqnarray*}
\int\sec xdx&=&\ln |\sec x+\tan x|+C\\
\int\csc x dx&=&\ln |\csc x-\cot x|+C\\
&=&-\ln |\csc x+\cot x|+C
\end{eqnarray*}

\noindent {\large \bf Trigonometric Identitties}

\begin{eqnarray*}
1&=&\sin^{2}\theta+\cos^{2}\theta\\
\sec^{2}\theta&=&1+\tan^{2}\theta\\
\sin (2\theta)&=&2\sin\theta\cos\theta\\
\cos\theta&=&\cos^{2}\theta-\sin^{2}\theta\\
\sin^{2}\theta&=&\frac{1}{2}(1-\cos 2\theta)\\
\cos^{2}\theta&=&\frac{1}{2}(1+\cos 2\theta)\\
\sin x\cos y&=&\frac{1}{2}(\sin (x-y)+\sin(x+y))\\
\sin x\sin y&=&\frac{1}{2}(\cos (x-y)-\cos(x+y))\\
\cos x\cos y&=&\frac{1}{2}(\cos (x-y)+\cos(x+y))\\
\sin(x\pm y)&=&\sin x\cos y\pm \cos x\sin y\\
\cos(x\pm y)&=&\cos x\cos y\mp \sin x\sin y\\
\end{eqnarray*}
 
\noindent {\large \bf Midpoint Rule}

\begin{eqnarray*}
\int_{a}^{b}f(x)dx&\simeq& M_{n}\\
&=&\Delta x (f(\overline{x}_{1})
+f(\overline{x}_{2})+\cdots +f(\overline{x}_{n}))
\end{eqnarray*}
$$(\Delta x=\frac{b-a}{n}, x_{i}=a+i\Delta x, \overline{x}_{i}
=\frac{x_{i-1}+x_{i}}{2})$$
Suppose $|f''(x)|\le K$ for $a\le x\le b$. If $E_{M}$ is 
the error in the middle point rule, then
$$|E_{M}|\le \frac{K(b-a)^{3}}{24n^{2}}.$$


\noindent {\large \bf Trapezoidal Rule}

\begin{eqnarray*}
\int_{a}^{b}f(x)dx&\simeq& T_{n}\\
&=&\frac{\Delta x}{2} 
(f(x_{1})
+2f(x_{2})+\cdots\\
&&\quad\quad\quad\quad
 +2f(x_{n-1})+f(x_{n}))
\end{eqnarray*}
$$(\Delta x=\frac{b-a}{n}, x_{i}=a+i\Delta x)$$
Suppose $|f''(x)|\le K$ for $a\le x\le b$. If $E_{T}$ is 
the error in the middle point rule, then
$$|E_{T}|\le \frac{K(b-a)^{3}}{12n^{2}}.$$

\noindent {\large \bf Simpson's Rule}

\begin{eqnarray*}
\int_{a}^{b}f(x)dx&\simeq& S_{n}\\
&=& \frac{\Delta x}{3} 
(f(x_{1})
+4f(x_{2})+2f(x_{3})+\cdots\\
&&\quad
 +2f(x_{n-2})+4f(x_{n-1})+f(x_{n}))
\end{eqnarray*}
$$(n \;\mbox{\rm even}, 
\Delta x=\frac{b-a}{n}, x_{i}=a+i\Delta x)$$
Suppose $|f^{(4)}(x)|\le K$ for $a\le x\le b$. If $E_{S}$ is 
the error in the middle point rule, then
$$|E_{S}|\le \frac{K(b-a)^{5}}{180n^{4}}.$$

\noindent {\large \bf Length of a Plane Curve}

The length 
of the curve $y=f(x)$ for $a\le x\le b$ is 
${\displaystyle \int_{a}^{b}\sqrt{1+(f'(x))^{2}}dx}$.

The length 
of the curve $x=g(y)$ for $c\le y\le d$ is 
${\displaystyle \int_{c}^{d}\sqrt{1+(g'(y))^{2}}dy}$.

The length of the parametric curve $x=x(t)$, $y=y(t)$ 
for $a\le t\le b$ is 
${\displaystyle \int_{a}^{b}\sqrt{\left(\frac{dx}{dt}\right)^{2}
+\left(\frac{dy}{dt}\right)^{2}}dt}$.

\noindent {\large \bf Polar coordinates}

$x=\cos \theta$, $y=\sin \theta$, $r=\sqrt{x^{2}+y^{2}}$,
$\theta=\tan^{-1}\left(\frac{y}{x}\right)$.

The area of the polar region 
$R=\{(r, \theta): f(\theta)\le r\le g(\theta), \alpha\le
\theta\le \beta\}$ is 
$$A=\frac{1}{2}\int_{\alpha}^{\beta}
((g(\theta))^{2}-(f(\theta))^{2})d\theta.$$


The length of a curve given by $r=f(\theta)$ for $\alpha\le
\theta\le \beta$ is ${\displaystyle \int_{\alpha}^{\beta}
\sqrt{r^{2}+\left(\frac{dr}{d\theta}\right)^{2}}d\theta}$.



\noindent {\large \bf Area of a Surface of Revolution}

Then the area
of the surface generated by rotating the curve 
$y=f(x)$ for $a\le x\le b$ about the $x$-axis is 
${\displaystyle \int_{a}^{b}2\pi f(x)\sqrt{1+(f'(x))^{2}}dx}$.

The area
of the surface generated by rotating the curve 
$x=g(y)$ for $c\le y\le d$ about the $y$-axis is 
${\displaystyle \int_{c}^{d}2\pi g(y)\sqrt{1+(g'(y))^{2}}dy}$.



\end{document}